% Section III --- paste into the CCNC manuscript.
% Citations use only keys from the current thebibliography.
% Point \ref{fig:e2e} at the existing end-to-end system diagram (Fig.~1).

\section{End-to-End: System Design}
\label{sec:system}

For trustworthy operational defense in a disaggregated enterprise network, we design a unified architecture with four static functional layers, as illustrated in Fig.~\ref{fig:e2e}: (1)~the \emph{Detection Layer}, where complementary observations at Access, Perimeter, and Endpoint are classified without pooling raw telemetry, and attribution names the evidence group that drove the decision; (2)~the \emph{Reasoning and Planning Layer}, where that detection product is turned into a policy-grounded, domain-specific mitigation plan; (3)~the \emph{Blockchain Governance Layer}, where the plan is cryptographically committed and the authorized action set is stored as an unpromptable whitelist; and (4)~the \emph{Execution and Monitoring Layer}, where an execution gate admits an action only if it is authorized for the predicted attack, identical to the committed plan, and integrity-checked against the on-chain digest.

Each network flow is processed in order through Detect, Reason, Commit, and Apply~\cite{li2025}. Reason does not authorize execution. Commit does not dispatch. Apply is the only layer that may touch a live control plane, and only after the gates succeed.

% III-A Detection --- paste-ready subsection.
% Citations use only keys from the current thebibliography.
% \ref{sec:reason} and \ref{fig:e2e} must exist in the manuscript.

\subsection{Detection Layer (Detect)}
\label{sec:detect}

What fails without this layer: a centralized detector requires each enterprise domain to surrender raw telemetry to a single point, which is exactly the constraint that makes cross-domain detection impractical to begin with. Regulatory, ownership, and trust boundaries prevent the Access, Perimeter, and Endpoint layers of an enterprise from pooling raw features in one place, so detection must happen without centralizing evidence, or it does not happen at all in a realistic multi-party setting. We address this with a three-party Vertical Federated Learning (VFL) architecture~\cite{liu2024vfl,folino2024,alazab2022}: each party keeps its own columns, a local multilayer perceptron (MLP) maps that slice to an embedding, and only the embedding reaches a coordinator. Raw values never leave the party that owns them.

The layer has a second, equally necessary job. A label alone is not actionable for Section~\ref{sec:reason}: knowing the event is a PortScan does not tell the planner which domain must act first. After the VFL classifier is trained, KernelSHAP is run on a distilled surrogate of the fused representation~\cite{wu2023falcon} so that $\hat{y}_i$ is attributed to one of the three evidence groups. The pair $(\hat{y}_i,\, p_i^{\star})$ is the only object forwarded to Reason.

To show that locality is not an accuracy tax, the same split, loss, and optimizer are used to train a centralized MLP on the concatenation of all features. That baseline never appears on the live path; it exists so the VFL-versus-pooled comparison isolates feature locality rather than a mismatch in model capacity~\cite{ye2024vfl}.
\begin{equation}
x_i=\{x_i^1,x_i^2,x_i^3\}
\;\xrightarrow{f_p}
h_i^p\in\mathbb{R}^{64}
\;\xrightarrow{\mathrm{concat}}
h_i\in\mathbb{R}^{192}
\;\xrightarrow{g}
\hat{y}_i
\;\xrightarrow{\mathrm{KernelSHAP}}
p_i^{\star}.
\label{eq:detect-pipeline}
\end{equation}

\paragraph{Feature Partitioning.}
Identifiers, addresses, and timestamps are dropped. The remaining 88 CICFlowMeter numeric columns from CICIDS-2017~\cite{sharafaldin2018} are assigned to three evidence groups by a rule-based classifier on feature names (Table~\ref{tab:detect-parties}). The catalog lists 97 assignments over 88 unique columns; nine signals are shared so one group can corroborate another without a raw-feature exchange. This is a statistical partition of a single flow-level feature set, not three independently instrumented systems. We present it under Access~/~Perimeter~/~Endpoint because those are the enterprise layers where each signal type would originate (Fig.~\ref{fig:e2e}). Original CICIDS labels are collapsed to $K=9$: BENIGN, DDOS, DOS, SSHPATATOR, FTPPATATOR, PORTSCAN, WEBATTACK, BOT, and OTHERS (classes with fewer than $200$ flows, including Infiltration and Heartbleed). Features at each party are standardized with a party-local scaler and stored as $32$-bit tensors. The aligned index set is split $64\%$~/~$16\%$~/~$20\%$ train~/~validation~/~test by stratified sampling (seed $42$), so every party sees the same flows and no party sees another party's columns.

\begin{table}[t]
\caption{VFL feature partition used as a controlled proxy for Access, Perimeter, and Endpoint ownership.}
\label{tab:detect-parties}
\centering
\begin{tabular}{clcl}
\hline
Party $p$ & Enterprise label & $d_p$ & Evidence \\
\hline
1 & Access~/~ISP & 23 & Volume/rate (duration, packet/byte totals, SYN counters) \\
2 & Perimeter~/~IDS & 32 & Packet-shape and ports (length, unique/vuln./HTTP/DNS) \\
3 & Endpoint~/~EDR & 42 & Timing/direction (IAT, bidirectional and reverse-path flags) \\
\hline
\end{tabular}
\end{table}

\paragraph{Local MLP Encoder.}
Each party $p\in\{1,2,3\}$ owns an encoder $f_p$ that never leaves that party. It is a two-layer MLP, not a linear projection:
\begin{equation}
x_i^p \in\mathbb{R}^{d_p}
\;\xrightarrow{W_1\in\mathbb{R}^{128\times d_p}}
\mathbb{R}^{128}
\;\xrightarrow{\mathrm{ReLU},\,\mathrm{Dropout}(0.2)}
\;\xrightarrow{W_2\in\mathbb{R}^{64\times 128}}
\mathbb{R}^{64}
\;\xrightarrow{\mathrm{ReLU},\,\mathrm{Dropout}(0.1)}
h_i^p.
\end{equation}
Only $h_i^p$ is transmitted. Three such encoders run in parallel with independent weights $\theta_p$. Concatenation---not summation---is used at the coordinator so each $64$-dimensional block remains identifiable for SHAP.

\paragraph{Coordinator MLP (VFL Classifier).}
The coordinator concatenates $h_i=[h_i^1\,\|\,h_i^2\,\|\,h_i^3]\in\mathbb{R}^{192}$ and applies a shared active classifier $g$:
\begin{equation}
h_i
\;\xrightarrow{192\rightarrow 128,\;\mathrm{ReLU},\;\mathrm{Dropout}(0.2)}
\;\xrightarrow{128\rightarrow K}
z_i,\qquad
\hat{y}_i=\arg\max_k z_{i,k}.
\end{equation}
Logits $z_i$ are unnormalized; softmax lives only in the loss and in the confidence $\mathrm{conf}_i=\max_k\mathrm{softmax}(z_i)_k$ that Reason will consume. The end-to-end VFL model is therefore three local MLPs plus one coordinator MLP, trained by back-propagating through the concatenation~\cite{liu2024vfl}. Gradients of $g$ return to each $f_p$; raw $x_i^p$ still never move.

\paragraph{Training.}
Both VFL and the centralized baseline use the same optimization so the comparison is not a training artifact. The loss is class-weighted cross-entropy. Inverse-frequency weights are clamped at $20\times$ so rare classes cannot dominate the gradient. The optimizer is Adam ($\mathrm{lr}=10^{-3}$, weight decay $10^{-4}$); gradients are clipped at $\ell_2=1$. A plateau scheduler halves the learning rate when validation loss stalls (factor $0.5$, patience $15$, floor $10^{-6}$). Training runs at most $100$ epochs, validates every $5$ epochs, and early-stops on macro-$F_1$ (patience $20$, $\Delta_{\min}=0.001$), restoring the best checkpoint. Reported metrics are accuracy, macro-recall, and macro-$F_1$, plus a per-class report against the centralized model.

\paragraph{Centralized All-Feature Baseline.}
The privacy-preserving claim is empty unless VFL is compared with a model that is allowed to see every column~\cite{ye2024vfl,yu2024vflprivacy}. We therefore concatenate the three already-scaled slices into one vector $\tilde{x}_i=[x_i^1\|x_i^2\|x_i^3]\in\mathbb{R}^{d}$ and train a non-federated MLP of comparable depth:
\begin{equation}
\tilde{x}_i
\;\rightarrow\; 256 \;\xrightarrow{\mathrm{ReLU},\,\mathrm{Dropout}(0.2)}\;
128 \;\xrightarrow{\mathrm{ReLU},\,\mathrm{Dropout}(0.2)}\;
64 \;\xrightarrow{\mathrm{ReLU},\,\mathrm{Dropout}(0.2)}\; K.
\end{equation}
It uses the same split, class weights, Adam settings, scheduler, gradient clip, and early-stopping rule. It is deliberately not a shallow network: a weaker centralized model would make VFL look better for the wrong reason. The centralized MLP is an evaluation instrument. It is not deployed, it does not produce SHAP for planning, and it does not feed Reason. The live path is VFL embeddings only.

\paragraph{Attribution via Distilled Surrogate.}
KernelSHAP on the live VFL stack would require hundreds of forward passes per explained flow, including three local encoders. That is impractical in a detection pipeline. We therefore freeze the trained VFL model as teacher, form the $192$-d concatenation of its embeddings, and train a compact surrogate $m$ ($192\rightarrow 128\rightarrow 64\rightarrow K$, dropout $0.2$ on the first hidden layer) by soft-label distillation: KL divergence between temperature-smoothed distributions at $T=3.0$, scaled by $T^2$, Adam as above, $50$ epochs. Teacher--student agreement (hard-label match, KL, logit MSE) is measured on the test set before any SHAP value is trusted~\cite{wu2023falcon}.

KernelSHAP then explains $m$, not the three encoders: background size $100$, perturbation budget $200$ per instance, up to $200$ test flows. Values lie in $\mathbb{R}^{192}$ and are summed inside each $64$-d block to one score per evidence group. The group with the largest mean $\lvert\phi\rvert$ is $p_i^{\star}\in\{\text{Access~/~ISP},\text{Perimeter~/~IDS},\text{Endpoint~/~EDR}\}$. Top-$3$ features per block and top-$5$ features overall are retained; those are the only feature names Section~\ref{sec:reason} is allowed to put in the retrieval query and the prompt.

\paragraph{Handoff to Reasoning.}
The Detection Layer emits, for each flow $i$,
\[
(\hat{y}_i,\; \mathrm{conf}_i,\; p_i^{\star},\; \text{three domain shares},\; \text{top features by }\lvert\phi\rvert).
\]
This closes the detection-without-centralization problem: Access, Perimeter, and Endpoint never pooled raw columns, and the coordinator still produced a class competitive with the all-feature MLP. The output is still evidence, not a plan~\cite{li2025}. Converting that record into a catalog action with a network tier---without inventing names or ignoring $p_i^{\star}$---is the failure the Reasoning and Planning Layer exists to close.


\subsection{Reasoning and Planning Layer (Reason)}
\label{sec:reason}

Section~\ref{sec:detect} closes detection-without-centralization and forwards a structured record, not a narrative: the predicted class $\hat{y}_i$, the detector confidence $\mathrm{conf}_i$, the dominant evidence group $p_i^{\star}$ obtained by summing KernelSHAP over each $64$-dimensional fusion block, the three domain shares, and the top local features by $\lvert\phi\rvert$. That record is still evidence. Knowing the event is a PortScan whose Perimeter block is largest does not name an authorized control, does not say whether Access or Endpoint should move second, and does not bind those choices to a catalog an execution gate can check. An unconstrained language model can invent a fluent mitigation that no playbook allows, or can tell a host-only story for a perimeter-dominated scan~\cite{xi2023,kshetri2025agentic}. We therefore treat Reasoning as a separate layer whose job is to convert the Detection product into one machine-checkable Mitigation Plans object. The layer does not authorize execution and does not write to the ledger; those failures belong to Commit and Apply.

The conversion is a fixed pipeline. After decode, nothing is free-form: the retrieval query is a filled template, ranking is a seven-step recall-then-precision cascade, the prompt is a slot-filled string, and the only accepted response is one JSON object whose action names must already appear in the per-attack catalog $W[\hat{y}_i]$ that Commit will later enforce on-chain~\cite{li2025}.
\begin{equation}
(\hat{y}_i,\, p_i^{\star},\, \mathrm{conf}_i,\, \mathrm{SHAP}_i)
\;\xrightarrow{\mathrm{query}}\;
q_i
\;\xrightarrow{\mathrm{retrieve+rank}}\;
\mathcal{R}_i
\;\xrightarrow{\mathrm{prompt}}\;
\pi_i
\;\xrightarrow{\mathrm{JSON}}\;
\mathrm{plan}_i.
\end{equation}

\paragraph{From Detection into a Query.}
Domain names in every later string are forced to exactly Access~/~ISP, Perimeter~/~IDS, or Endpoint~/~EDR---the same three evidence groups used to split the 88 CICFlowMeter columns and to aggregate SHAP (Fig.~\ref{fig:e2e}). $\hat{y}_i$ selects the allowed-action list $W[\hat{y}_i]$. $\mathrm{conf}_i$ is written into both the query and the task block so intensity (\texttt{Immediate} $\ldots$ \texttt{Low}) tracks the detector rather than the model's tone. $p_i^{\star}$ becomes the sentence ``Dominant network domain: $\ldots$ (contribution: $x\%$)''; top-3 features per domain become retrieval keywords, and only the top-5 features by $\lvert\phi\rvert$ are admitted into the evidence JSON the model may cite. Retrieval does not ask the language model to invent a search string. The operational query is a single template filled from that record. Optional rephrases exist; the evaluated path uses only the template, so two identical Detection records produce the same $q_i$. That determinism is required by the next layer: Commit hashes a plan that must be reproducible from the same evidence.

\paragraph{Policy Index.}
The corpus is closed---NIST SP~800-53~\cite{nist80053}, NIST SP~800-207~\cite{nist800207}, CIS Controls~\cite{cisv8}, and MITRE ATT\&CK~\cite{mitre_attack}---not an open web~\cite{shajarian2026,fu2026}. Documents are stored parent/child. A parent is a semantic section (the text the planner will read). A child is a short passage used only for nearest-neighbor search (size $384$, overlap $96$, embedding \texttt{all-MiniLM-L6-v2}), tagged with $(\mathrm{parent\_id},\,\mathrm{child\_index},\,\mathrm{source})$. Ranking operates on children; the prompt receives expanded parents.

\paragraph{Retrieve-then-Rank Cascade.}
Dense search alone is the wrong ranking: a bi-encoder returns topical neighbors and, left unchecked, fills the prompt with near-duplicates from one source. Each step closes a failure of the previous one. (1)~\emph{Dense retrieve:} FAISS oversamples ($\min(300,\,20\times 6)$), keeps $20$ children, converts distance $d$ to $1/(1+d)$, and balances across source files. (2)~\emph{Merge and dedupe} by parent/child identity so overlap does not starve the later budget. (3)~\emph{MMR} with $\lambda=0.5$ keeps up to $60$ children, $\mathrm{MMR}(c)=\lambda\cdot\mathrm{sim}(q,c)-(1-\lambda)\cdot\max_{c'\in S}\mathrm{sim}(c,c')$, so NIST, CIS, and ATT\&CK can coexist. (4)~\emph{Cross-encoder} \texttt{ms-marco-MiniLM-L-6-v2} scores each $(q,\mathrm{passage})$ jointly (passage truncated at $4000$ characters)---vector similarity is a neighborhood, not ``this passage mitigates \emph{this} PortScan at the perimeter.'' It is too expensive for the whole index, which is why (1)--(3) exist. (5)~Keep the top $20$ children. (6)~Expand each to its parent ($12\,000$-character cap); duplicate parents collapse. (7)~Emit the top five numbered sections, or the explicit empty-KB sentence. Retrieval returns guidance; it does not invent action names. $W[\hat{y}_i]$ is a different prompt slot and the same catalog Commit will store on-chain.

\paragraph{Prompt Construction.}
The user message is one template filled in software before decode~\cite{xi2023,li2025}. The model never sees raw telemetry, never sees the $192$-d fusion vector, and never sees names outside $W[\hat{y}_i]$. Slots are: role and the three exact domain strings; prediction summary; the dominance sentence from $p_i^{\star}$; Detection evidence JSON (top-5 features only); $W[\hat{y}_i]$ plus preferred domains and evidence cues; a per-domain card (role, top-3 features, actions that domain may execute); numbered RAG sections or the empty-KB sentence. The task block requires interpretation on all three domains, a first-mover domain, catalog-only actions, a network tier on each unit, evidence-tied reasoning, priority scaled by $\mathrm{conf}_i$, and empty lists rather than a made-up control. Decoding uses GPT-4o-mini at $T=0.3$; the first ``\{'' through the last ``\}'' is parsed.

\paragraph{Response Schema.}
The only accepted output is $\mathrm{plan}_i$: $\mathtt{threat\_level}\in\{\mathrm{Critical},\mathrm{High},\mathrm{Medium},\mathrm{Low}\}$; primary and supporting units of the form $\{\mathtt{action},\,\mathtt{network\_tier},\,\mathtt{party\_evidence\_type},\,\mathtt{reasoning}\}$; $\mathtt{all\_actions}$ their union; $\mathtt{overall\_reasoning}$; $\mathtt{execution\_priority}\in\{\mathrm{Immediate},\mathrm{High},\mathrm{Standard},\mathrm{Low}\}$; and $\mathtt{knowledge\_sources\_used}$. The field $\mathtt{action}$ is copied from $W[\hat{y}_i]$ (\texttt{limit rate} remains \texttt{limit rate}). $\mathtt{network\_tier}$ is exactly one of the three domain strings. Primary actions prefer $p_i^{\star}$ unless the evidence contradicts it. A PortScan whose Perimeter share dominates is planned as, for example, \texttt{tarpit scan} and \texttt{harden ports} at Perimeter~/~IDS with supporting \texttt{update ACL} at Access~/~ISP---names legal for PORTSCAN and illegal for BENIGN. Apply will reject any name absent from $W[\hat{y}_i]$ or absent from this object.

\paragraph{Handoff to Blockchain Governance.}
$\mathrm{plan}_i$ is this layer's only output. It is the off-chain payload $P$ that Section~\ref{sec:commit} will canonicalize and commit as $c=\mathrm{SHA256}(P)$. Reason never calls the contract and never opens a device API. This closes the evidence-to-plan failure that Section~\ref{sec:detect} left open. It does not close the trust failure that follows generation. A fluent, catalog-legal action copied from another cycle, a rewritten network tier, or a payload whose bytes change before dispatch will still look like a valid plan if the next layer only stores text. Binding $W[\hat{y}_i]$ on-chain, writing $c$ once, and refusing any apply that does not re-hash to $c$ is the failure the Blockchain Governance Layer exists to close.

% III-C Blockchain Governance --- paste-ready subsection.
% Citations use only keys from the current thebibliography.
% \ref{sec:reason}, \ref{sec:apply}, \ref{fig:e2e} must exist in the manuscript.

\subsection{Blockchain Governance Layer (Commit)}
\label{sec:commit}

Section~\ref{sec:reason} forwards a single object $\mathrm{plan}_i$: the Mitigation Plans JSON whose action names were copied from $W[\hat{y}_i]$ and whose network tier is exactly Access~/~ISP, Perimeter~/~IDS, or Endpoint~/~EDR. That object is still a \emph{proposal}. Reasoning is probabilistic. A language model can emit a fluent action that is still harmful on a live path; a poisoned or substituted plan can change between generation and dispatch; a policy paragraph in a prompt can be ignored~\cite{xi2023,lazer2026,chhabra2026}. Autonomous defense therefore needs a pre-execution governance layer that cannot hallucinate mitigations and cannot be talked into accepting them~\cite{karim2025,pendli2025,jan2025}.

A smart contract is a program that lives on a ledger and executes exactly the rules written in its code. Callers submit transactions; the network orders them; every honest replica applies the same function to the same state. We use that property for two jobs only: an unpromptable per-attack whitelist, and a write-once digest of the plan produced for this detection cycle~\cite{zhang2024exploit,chen2024,jin2024}. The contract does not classify traffic, does not retrieve policy, and does not dispatch. Those remain Detect, Reason, and Apply.
\begin{equation}
\underbrace{\mathrm{plan}_i}_{\text{III-B}}
\;\xrightarrow{\mathrm{canonicalize}}\;
P
\;\xrightarrow{\mathrm{SHA256}}\;
c
\;\xrightarrow{\mathrm{anchor}}\;
\mathrm{ledger}[k_{\mathrm{agent}},k_{\mathrm{report}}]
\;\xrightarrow{\mathrm{TrustAnchored}}\;
\text{SOC / Apply}.
\label{eq:commit-pipeline}
\end{equation}

\paragraph{Off-Chain Payload versus On-Chain Digest.}
After Reason returns, the Mitigation Plans object is persisted off-chain (identity, structured plan, retrieval context). The backend then builds a canonical payload $P$ with fixed keys, sorted lexicographically, compact separators, and UTC timestamps so two encodings of the same report cannot produce two hashes: payload version; report and prediction identities; optional row index (the same Detection row Reason used); creation time; structured $\mathrm{plan}_i$ (threat level, primary and supporting units, tiers, reasoning); and the retrieval context used in the prompt.
\begin{equation}
c = \mathrm{SHA256}(\mathrm{canonicalJSON}(P)).
\label{eq:commitment}
\end{equation}
Telemetry, the $192$-d fusion vector, SHAP tables, and free-form model prose are not required on the ledger~\cite{zhang2024exploit}. Only the $32$-byte digest $c$ is written. Confidentiality is the point: an auditor who holds $P$ can re-hash and compare; an observer of the chain cannot reconstruct the flow or the prompt.

\paragraph{Addressing and Keys.}
The registry never stores human-readable identifiers. Four $\mathtt{bytes32}$ keys are derived before the transaction (Table~\ref{tab:commit-keys}). Lookups are $W[k_a][k_u]$ and $\mathrm{commitments}[k_{\mathrm{agent}}][k_{\mathrm{report}}]$. A paraphrase that Reason was forbidden to emit also fails here: \texttt{limit rate} and \texttt{rate limiting} are different keys.

\begin{table}[t]
\caption{Hash keys used by the governance registry. Attack and action tokens match the UTF-8 labels used by Detect and Reason.}
\label{tab:commit-keys}
\centering
\begin{tabular}{p{0.18\columnwidth}p{0.36\columnwidth}p{0.36\columnwidth}}
\hline
Key & Formation & Role \\
\hline
$k_{\mathrm{agent}}$ & SHA-256 of the agent-job identity & Bind the cycle to the requester \\
$k_{\mathrm{report}}$ & SHA-256 of the report identity & One commitment slot per plan \\
$k_a$ & $\mathrm{keccak256}(\hat{y}_i)$ & Attack class as the contract sees it \\
$k_u$ & $\mathrm{keccak256}(u)$ & Exact catalog token (not a paraphrase) \\
\hline
\end{tabular}
\end{table}

\paragraph{Job 1: Per-Attack Whitelist.}
At deploy, the owner seeds $W[a,u]$ from the same catalog Reason was allowed to name (CICIDS-2017 grouping~\cite{sharafaldin2018}). After seed, ordinary callers cannot add a name. $\mathtt{isActionWhitelisted}(k_a,k_u)$ is a view; it does not depend on the prompt. An action that is legal for DDoS and illegal for the predicted PortScan fails even if the sentence is fluent. That is the failure a policy paragraph cannot close: the model cannot talk a new token onto $W$~\cite{jan2025,singh2025zt}.

\paragraph{Job 2: Write-Once Anchor.}
$\mathtt{anchor}(k_{\mathrm{agent}},k_{\mathrm{report}},c)$ requires every argument nonzero and $\mathrm{commitments}[k_{\mathrm{agent}}][k_{\mathrm{report}}]=0$. A second write is rejected. The contract emits $\mathtt{TrustAnchored}$. Later readers call $\mathtt{getCommitment}$ and compare to a fresh SHA-256 of the stored $P$. Three outcomes are distinguishable: (i)~hashes match---$P$ is the plan that was committed; (ii)~hashes differ---$P$ was edited after Commit; (iii)~commitment is empty---this cycle was never anchored, so Apply must not proceed.

Immutability here is not ``the blockchain stores the network.'' It is three enforceable facts: the allowed action set for each attack is on-chain and not promptable; the plan for this cycle has a single committed digest; any later apply must bind to both~\cite{alqithami2026,pendli2025}.

\paragraph{What the Contract Will Not Check.}
$\mathtt{applyAction}(k_a,k_u,k_{\mathrm{agent}},k_{\mathrm{report}})$ is a receipt, not the full gate. On-chain it requires (i)~$W[k_a,k_u]=1$, (ii)~a nonzero commitment for that pair, and (iii)~the unit not already applied. It does not restore $\tau$ or test $u\in\mathrm{plan}(P)$. Those comparisons need the off-chain JSON. Software therefore evaluates, in order, whitelist, plan-binding ($u$ and $\tau$ identical to Reason's unit), and digest ($c_{\mathrm{chain}}=c_{\mathrm{rehash}}$) \emph{before} the receipt transaction. A still-whitelisted action copied from another cycle dies at plan-binding; a rewritten payload dies at the digest. Putting the full JSON on-chain would break the hash-only confidentiality claim. The split is intentional: the ledger freezes $W$ and $c$; Apply restores $\mathrm{plan}_i$ (Fig.~\ref{fig:e2e}).

\paragraph{Governance as a Signal, Not an Actuator.}
$\mathtt{TrustAnchored}$ is a tamper-evident event that a validated proposal exists. The SOC~/~SIEM~/~monitoring bus is a subscriber, not a vendor-specific product~\cite{chen2024,jin2024}. Review/Approve from an operator is recorded back through the same registry and re-enters retrieve--compare; approval does not skip the gate. The contract freezes what may be done. It does not open API, gRPC, or MCP.

\paragraph{Handoff to Execution and Monitoring.}
Commit emits $c$, the event, and the keys Section~\ref{sec:apply} will need to retrieve $W[\hat{y}_i]$ and $\mathtt{getCommitment}$. This closes the ungoverned-proposal failure that Section~\ref{sec:reason} left open: $\mathrm{plan}_i$ can no longer be silently replaced by another fluent JSON. It does not close actuation. A still-whitelisted unit aimed at the wrong domain, or a second apply of the same $(k_{\mathrm{agent}},k_{\mathrm{report}},k_u)$, must be refused at the execution gate. Restoring $\tau$, testing $u\in\mathrm{plan}(P)$, and only then calling $\mathtt{applyAction}$ is the failure the Execution and Monitoring Layer exists to close.


\subsection{Execution and Monitoring Layer (Apply)}
\label{sec:apply}

What fails without this layer: Commit records that a plan existed, but recording is not enforcement. Without an execution gate, an orchestrator could still dispatch a different legal action, a silently edited payload, or a unit aimed at the wrong enterprise domain. Apply is the only place a planned action may touch the live network. After the SOC or core monitoring path is notified, an event listener retrieves the on-chain record, the gate compares that record with policy, and only then may the orchestrator dispatch (Fig.~\ref{fig:e2e}).

\paragraph{Execution Gate.}
The gate does not invent mitigations. It compares three facts that already exist: (i)~whether action $u$ is on the whitelist $W$ for this attack; (ii)~whether $u$ and its domain $\tau$ (Access, Perimeter, or Endpoint) are the same actions Reason wrote into plan $P$; (iii)~whether a re-hash of the off-chain payload still matches the on-chain digest. In compact form,
\begin{equation}
W[a,u]=1 \;\wedge\; u\in\mathrm{plan}(P) \;\wedge\; c_{\mathrm{chain}}=c_{\mathrm{rehash}}.
\label{eq:apply-gate}
\end{equation}
Software evaluates the three checks in that order---whitelist, then plan-binding, then digest---before the registry is asked to record an apply. The on-chain apply then re-checks that the action is still whitelisted, that a commitment exists, and that the same unit has not already been applied. An action that is only plausible, or that is legal for a different attack, fails the whitelist. A still-whitelisted action copied from another cycle fails plan-binding. A silently edited payload fails the digest check.

\paragraph{Per-Attack Whitelist.}
Labels follow the CICIDS-2017 grouping used by Detect~\cite{sharafaldin2018}. Actions are a constrained catalog, not free-form model text (Table~\ref{tab:whitelist}).

\begin{table}[t]
\caption{Per-attack action whitelist $W$ at Apply time. Labels follow CICIDS-2017~\cite{sharafaldin2018}; actions are a constrained catalog, not free-form LLM text.}
\label{tab:whitelist}
\centering
\begin{tabular}{p{0.28\columnwidth}p{0.64\columnwidth}}
\hline
Attack $a$ & Allowed actions (excerpt) \\
\hline
BENIGN & log incident, monitor traffic \\
DDOS~/~DOS & limit rate, enable scrubbing, block IP, enable syncookies, update ACL \\
SSHPATATOR~/~FTPPATATOR & fail2ban block, lock account, enforce MFA \\
PORTSCAN & tarpit scan, harden ports, update ACL \\
WEBATTACK & apply WAF, virtual patch, isolate service \\
BOT & captcha challenge, reputation filter \\
OTHERS & log incident, isolate service, update ACL \\
\hline
\end{tabular}
\end{table}

\paragraph{Reject Path and Human Loop.}
If the gate does not match, the plan is marked rejected or blocked. The flow goes to operator review, not to the network. Every such outcome is written to the same monitoring archive as a successful apply, so a denial is as visible as an actuation. If an operator later reviews and approves, that approval is recorded back through the governance layer and re-enters the same retrieve--compare path. Approval does not skip the gate.

\paragraph{Success Path.}
If the gate validates the plan, control passes to the execution agent. That agent invokes the domain-specific interface---API, gRPC, or MCP---for the Access, Perimeter, or Endpoint tier named in the plan. A successful unit is recorded on-chain and refused if presented again. The evaluated executor records receipts rather than driving a live device; a production deployment can replace that stub with the same control-plane fabric that collected the flow. Dispatch is therefore a signed, gated call, not a free-form command bus.

\begin{algorithm}[t]
\caption{Commit and Apply}
\label{alg:commit-apply}
\begin{algorithmic}[1]
\PROCEDURE{Commit}{$\mathit{agentKey}$, $\mathit{reportKey}$, $c$}
    \STATE \textbf{require} $c \neq 0$ and not already committed
    \STATE store commitment $c$
    \STATE emit trust-anchored
\ENDPROCEDURE
\PROCEDURE{Apply}{$a$, $u$, $\tau$, $\mathit{agentKey}$, $\mathit{reportKey}$}
    \STATE \textbf{require} $\tau$ matches the domain written in $\mathrm{plan}(P)$
    \STATE \textbf{require} $W[a,u]=1$ \COMMENT{policy ground}
    \STATE \textbf{require} a commitment exists and this unit is not already applied
    \STATE \textbf{require} $u \in \mathrm{plan}(P)$ and $c_{\mathrm{chain}}=\mathrm{SHA256}(P)$
    \STATE mark applied; emit action-applied
    \STATE dispatch $u$ to domain $\tau$
\ENDPROCEDURE
\end{algorithmic}
\end{algorithm}

Every hop---notify, retrieval, whitelist miss, plan mismatch, human approval, successful apply---emits an audit record back to the monitoring core. The core archive is the operator-visible history; the ledger holds the immutable receipts. This layer guarantees that only a whitelisted, bound, integrity-checked action reaches an agent. It does not claim that the language-model plan was correct.
